\scanpage{7}
\uline{Généralisation sur une base quelconque.}

Une \uline{courbe standart} sur $S$ (multiplicité schématique, disons) sera défini constructivement en termes d'un système a) b) c) \inserted{comme} ci-dessus, i.e.
\[
(5)\quad\left\{\begin{aligned}
&\text{a) }Y,\text{ schéma propre et lisse }\inserted{\text{sur }S}\text{ partout de}\\
&\text{dim relative }1;\\
&\text{b) }\widetilde A\text{ et }I,\text{ deux sous-schém. de }Y,\\
&\text{finis étales sur }S,\quad\deleted{\unsure{\text{avec}}}\ \widetilde A\cap I=\varnothing;\\
&\text{c) }\sigma_{\widetilde A}\text{ une involution de }\widetilde A,\text{ opérant}\\
&\text{librement.}
\end{aligned}\right.
\]
On construit alors $\widehat X=Y/\sigma$, contenant $A\simeq\widetilde A/\sigma$ et $I$ comme sous-schémas \inserted{fermés} finis étales sur $S$, et on définit $X=\widehat X\setminus I$. On peut montrer que le foncteur
\[
(6)\qquad(Y,\widetilde A,\inserted{I,}\sigma_{\widetilde A})\longmapsto X
\]
\margin{\uline{NB} NB ce foncteur est compatible \unsure{aux} images inverses\ldots}
des systèmes (5) (pour les iso) vers les schémas relatifs, \inserted{(pour les iso)} est \uline{pleinement fidèle}.
\margin{faux tel quel !}
\mspage{6}
